\documentclass[12pt,a4paper]{article}

\usepackage{graphicx}
\usepackage[a4paper,margin=0.8in]{geometry}
\usepackage{changepage}
\usepackage{float}
\usepackage{url}
\usepackage{amsmath}
\usepackage{tabularx}
\usepackage{titlesec}
\usepackage{listings}
\usepackage{xcolor}

\definecolor{codegreen}{rgb}{0,0.6,0}
\definecolor{codegray}{rgb}{0.5,0.5,0.5}
\definecolor{codepurple}{rgb}{0.58,0,0.82}
\definecolor{backcolour}{rgb}{0.95,0.95,0.92}

\lstdefinestyle{mystyle}{
    backgroundcolor=\color{backcolour},   
    commentstyle=\color{codegreen},
    keywordstyle=\color{magenta},
    numberstyle=\tiny\color{codegray},
    stringstyle=\color{codepurple},
    basicstyle=\ttfamily\footnotesize,
    breakatwhitespace=false,         
    breaklines=true,                 
    captionpos=b,                    
    keepspaces=true,                 
    numbers=left,                    
    numbersep=5pt,                  
    showspaces=false,                
    showstringspaces=false,
    showtabs=false,                  
    tabsize=4
}
\lstset{style=mystyle}

\titleformat{\section}
  {\bfseries}
  {\thesection.}{0.5em}{}

\titleformat{\subsection}
  {\large\bfseries}
  {\thesubsection}{1em}{}


\title{Computer Architecture Lab 04}
\author{K.D.N.Chandrasena}
\date{\today}

\begin{document}

\begin{center}
{\Large \textbf{Department of Computer Engineering}}\\
\vspace{0.2cm}
{\large University of Peradeniya}\\
\vspace{0.6cm}

{\large CO2070 -- COMPUTER ARCHITECTURE}\\
\vspace{0.3cm}
{\Large \textbf{LAB 04 REPORT}}\\
\vspace{0.2cm}

\end{center}
{\large
\textbf{Team Number:} No. 15\\[0.3cm]
\textbf{Name:} K.D.N.Chandrasena , Tharuka Cooray\\[0.3cm]
\textbf{Registration Number:} E/22/054 , E/2//058\\[0.3cm]
\textbf{Date:} \today
}

\vspace{1cm}

% ============================================================
% SECTION 1
% ============================================================
\section*{Lab 04: Flow Control Instructions}

In this lab, the processor datapath and control unit were upgraded to support flow control instructions, breaking the sequential execution of code. The two implemented instructions were the unconditional Jump (\texttt{j}) and the Branch If Equal (\texttt{beq}).

\subsection*{1. ALU Zero Flag Modification}
To support conditional branching, the ALU was modified to output a 1-bit \texttt{ZERO} flag. During the execution of a \texttt{beq} instruction, the control unit forces the datapath to subtract the two source registers. The \texttt{ZERO} flag continuously monitors the output.

\begin{lstlisting}[language=Verilog, caption=ALU Zero Flag Logic]
// If the ALU result is exactly zero, output 1
assign ZERO = (RESULT == 8'b00000000) ? 1'b1 : 1'b0;
\end{lstlisting}

\subsection*{2. Branch/Jump Target Adder}
A dedicated hardware adder was integrated into the datapath to calculate target addresses. The jump offset specifies the distance in instructions. The hardware shifts this offset left by 2 bits to convert it to a byte distance and adds it to the \texttt{PC+4} value.

\begin{lstlisting}[language=Verilog, caption=Target Address Adder]
// Base address calculation
assign #1 pc_next = PC + 32'h4;

// Target calculation: Shift offset (bits 23:16) left by 2 and sign-extend
assign #2 pc_target = pc_next + {{22{INSTRUCTION[23]}}, INSTRUCTION[23:16], 2'b00};
\end{lstlisting}

\subsection*{3. Program Counter Multiplexer}
A multiplexer was added before the Program Counter register. This MUX evaluates the control signals from the Control Unit and the \texttt{ZERO} flag from the ALU to determine the next execution address.

\begin{lstlisting}[language=Verilog, caption=PC Source Multiplexer]
// Select target address if Jump OR (Branch AND Zero)
assign pc_next_final = (jump | (branch & zero)) ? pc_target : pc_next;
\end{lstlisting}

\subsection*{4. Control Unit Decoding}
The Control Unit was updated to decode the opcodes for \texttt{j} and \texttt{beq}, routing the appropriate electrical signals into the datapath to trigger the multiplexers and ALU modes.

\begin{lstlisting}[language=Verilog, caption=Control Unit Opcode Decoding]
case (opcode)
    8'h06: begin                // j instruction
        JUMP = 1;
    end
    8'h07: begin                // beq instruction
        BRANCH = 1;             
        ALUOP = 3'b001;         // Set ALU to Addition
        Op2Select = 2'b01;      // Select negated OUT2 for subtraction
    end
endcase
\end{lstlisting}

\end{document}
